% GENERATED FILE. Edit canonical lessons, then run npm run generate:books.
% canonical-source-hash: fnv1a64:616850b3d7312727
% canonical-lessons: LA-C281-vidi-vidit, LA-C281-eram, LA-C281-ambulabam, LA-C281-pluebat-venit, LA-R281-the-scene-and-the-event

\chapter{The Scene and the Event}
\label{ch:la-the-scene-and-the-event}
\hlchaptermodality{281}

\begin{chapteropening}
\textbf{By the end of this chapter:} \emph{I can tell what happened in the perfect, and paint the scene it happened in with the imperfect.}

Complete LA-R281-the-scene-and-the-event: I can tell what happened in the perfect, and paint the scene it happened in with the imperfect.
\end{chapteropening}

\section[vīdī, vīdit, vīdērunt]{vīdī, vīdit — the past of what happened}

\label{lesson:LA-C281-vidi-vidit}



\begin{quote}
Say \textquotedblleft{}I recently read a book\textquotedblright{}. (\emph{Nūper librum lēgī.}) Now \textquotedblleft{}before, I lived in Rome\textquotedblright{}. (\emph{Anteā Rōmae habitāvī.}) And ask \textquotedblleft{}what happened?\textquotedblright{} (\emph{Quid accidit?})

All three are the \textbf{perfect}, Latin's past for a thing that happened and was over. Here are its six persons.
\end{quote}

\begin{grammarlens}[title={Grammar Lens: one stem, six endings}]
Every verb has a \textbf{perfect stem}: for \emph{videō}, \textbf{vīd-}, with a long \emph{ī}. Onto it go six endings, the same for every verb:

\begin{quote}
\textbf{vīdī} — I saw. \textbf{vīdistī} — you saw. \textbf{vīdit} — he saw, she saw. \textbf{vīdimus} — we saw. \textbf{vīdistis} — you all saw. \textbf{vīdērunt} — they saw.
\end{quote}

Stories lean on two of them: \textbf{-it} for one person, \textbf{-ērunt} for several.

The stem is learned with each verb. \emph{veniō} has \textbf{vēn-}: \emph{vēnit}, \emph{vēnērunt}. \emph{sum} has \textbf{fu-}: \emph{fuī}, I was; \emph{fuit}, he was.

\emph{Vīdī} is also \textquotedblleft{}I have seen\textquotedblright{}; the sentence decides.
\end{grammarlens}

\subsection*{You'll want to know: the -āvī verbs}
Most \textbf{-āre} verbs add \textbf{-āv-}, like the \emph{habitāvī} you own:

\begin{quote}
ambulāre — \textbf{ambulāvī}, I walked; \textbf{ambulāvit}, he walked clāmāre — \textbf{clāmāvit}, he shouted nārrāre — \textbf{nārrāvit}, she told; \textbf{nārrāvērunt}, they told
\end{quote}

The stems learned one by one are the others: \emph{vīdī}, \emph{vēnī}, \emph{fuī}, \emph{lēgī}, and \emph{bibī}, from \emph{bibō}, with nothing added at all.

\begin{cousinweb}
\textbf{fuī} looks nothing like \emph{sum} because it is built on another root, the one the chapter on \emph{sum} named: Proto-Indo-European \emph{*b\textsuperscript{h}uH-}, \textquotedblleft{}to become\textquotedblright{}, which EtymDB 2.1 lists among the parents of \emph{sum} and which also gives English \emph{be}.

\textbf{vīdī} keeps the root of \emph{videō} with a longer vowel (Lewis \& Short give \emph{vīdi} as its perfect). That vowel is how the ear tells \emph{videt}, \textquotedblleft{}he sees\textquotedblright{}, from \emph{vīdit}, \textquotedblleft{}he saw\textquotedblright{}.
\end{cousinweb}

\subsection*{Your turn}
Say these aloud:

\begin{itemize}
\raggedright
  \item \emph{vīdī, vīdistī, vīdit} — then \emph{vīdimus, vīdistis, vīdērunt}
  \item \emph{videt}, he sees; then \emph{vīdit}, he saw — and hear the long \emph{ī}
  \item \emph{vēnit}, he came; \emph{vēnērunt}, they came
  \item \emph{fuī}, I was; \emph{fuit}, it was
  \item \emph{ambulāvī}, I walked; \emph{ambulāvērunt}, they walked
  \item \emph{Nūper amīcum vīdī.} — Recently I saw a friend
\end{itemize}

\begin{tcolorbox}[breakable,colback=teal!4,colframe=teal!35!black,title={Before you move on}]
\textquotedblleft{}He saw\textquotedblright{}? (\emph{\textbf{Vīdit.}}) \textquotedblleft{}They came\textquotedblright{}? (\emph{\textbf{Vēnērunt.}}) \textquotedblleft{}I walked\textquotedblright{}? (\emph{\textbf{Ambulāvī.}}) Which ending says \textquotedblleft{}they\textquotedblright{}? (\emph{\textbf{-ērunt.}}) And which root is behind \emph{fuī}? (\emph{\textbf{*b\textsuperscript{h}uH-}}, \textquotedblleft{}to become\textquotedblright{}, the root of English \emph{be}.)
\end{tcolorbox}
\section[eram, erat, erant]{eram, erat — the past that stays open}

\label{lesson:LA-C281-eram}



\begin{quote}
Say \textquotedblleft{}I am\textquotedblright{} and \textquotedblleft{}they are\textquotedblright{}. (\emph{Sum. Sunt.}) Now \textquotedblleft{}I was\textquotedblright{}, in the perfect. (\emph{Fuī.})

\emph{Fuī} closes a door: I was, and that is over. Latin has a second \textquotedblleft{}was\textquotedblright{} that holds it open.
\end{quote}

\begin{grammarlens}[title={Grammar Lens: six forms, and the endings you know}]
\begin{quote}
\textbf{Ōlim Rōma parva erat.} — Long ago Rome was small.
\end{quote}

\textbf{erat} is \textquotedblleft{}was\textquotedblright{}, the \textbf{imperfect} of \emph{sum}. All six persons:

\begin{quote}
\textbf{eram} — I was. \textbf{erās} — you were. \textbf{erat} — he, she, it was. \textbf{erāmus} — we were. \textbf{erātis} — you all were. \textbf{erant} — they were.
\end{quote}

The endings \textbf{-m}, \textbf{-s}, \textbf{-t}, \textbf{-mus}, \textbf{-tis}, \textbf{-nt} are those of \emph{sum}, \emph{es}, \emph{est}, \emph{sumus}, \emph{estis}, \emph{sunt}, on the stem \textbf{era-}. Before \emph{-m}, \emph{-t} and \emph{-nt} the \emph{a} is short.

\emph{Erat} is a state that went on, with no edge in view.

\begin{quote}
\textbf{Fessus eram.} — I was tired. \textbf{Caelum obscūrum erat.} — The sky was dark.
\end{quote}
\end{grammarlens}

\begin{culture}
Both are \textquotedblleft{}was\textquotedblright{}; the choice is the frame. \emph{Laetus eram} paints a mood that lasted through the story. \emph{Laetus fuī} sums it up from outside: I was happy, and that time is over. A story sets its scene with \emph{erat}.
\end{culture}

\begin{cousinweb}
\textbf{eram} is usually explained as an older \emph{*esam}, built on the same root as \emph{est}: between two vowels an old Latin \emph{s} turned into \emph{r}, and \emph{es-a-m} came out as \emph{eram}. Neither EtymDB nor Lewis \& Short records a form that old, so this is the account histories of Latin give, offered here as one. The \emph{s} survives where no vowel followed it: \emph{es}, \emph{est}, \emph{estis}.
\end{cousinweb}

\subsection*{Your turn}
Say these aloud:

\begin{itemize}
\raggedright
  \item \emph{eram, erās, erat} — then \emph{erāmus, erātis, erant}
  \item \emph{Ōlim Rōma parva erat.}
  \item \emph{Fessus eram.} — then a woman says it: \emph{Fessa eram.}
  \item \emph{Caelum obscūrum erat.} — then \emph{In vīllā erāmus.}
  \item the closed door, then the open one — \emph{fuit}, then \emph{erat}
\end{itemize}

\begin{tcolorbox}[breakable,colback=teal!4,colframe=teal!35!black,title={Before you move on}]
\textquotedblleft{}They were\textquotedblright{}? (\emph{\textbf{Erant.}}) \textquotedblleft{}We were\textquotedblright{}? (\emph{\textbf{Erāmus.}}) Whose endings does \emph{eram} borrow? (The present \emph{\textbf{sum, es, est}}.) Which \textquotedblleft{}was\textquotedblright{} paints a scene, \emph{fuit} or \emph{erat}? (\emph{\textbf{Erat.}}) And what did the \emph{r} of \emph{eram} used to be? (An \emph{\textbf{s}}, between two vowels.)
\end{tcolorbox}
\section[ambulābam, habēbam, legēbam]{ambulābam — I was walking}

\label{lesson:LA-C281-ambulabam}



\begin{quote}
Say \textquotedblleft{}I was tired\textquotedblright{}, the open past. (\emph{Fessus eram.}) Now \textquotedblleft{}I walked\textquotedblright{}, the closed one. (\emph{Ambulāvī.})

\emph{Eram} belongs to one verb. Every other verb makes its open past with two letters slipped in before the ending.
\end{quote}

\begin{grammarlens}[title={Grammar Lens: -ba- before the ending}]
Take the present stem, add \textbf{-ba-}, then the endings \emph{eram} taught you:

\begin{quote}
\textbf{ambulābam} — I was walking. \textbf{ambulābās} — you were walking. \textbf{ambulābat} — he was walking. \textbf{ambulābāmus} — we were walking. \textbf{ambulābātis} — you all were walking. \textbf{ambulābant} — they were walking.
\end{quote}

The vowel before \emph{-ba-} comes from the infinitive. An \textbf{-āre} verb keeps its \emph{ā}, and an \textbf{-ēre} verb its \emph{ē}: \emph{habēre}, \textbf{habēbam}, I had, I used to have; \emph{sedēre}, \textbf{sedēbat}, he was sitting.

The other verbs put \textbf{-ēbam} on the stem: \emph{legere}, \textbf{legēbam}, I was reading; \emph{volō} gives \textbf{volēbat}, he wanted. And the verbs with \textbf{-iō} keep their \emph{i}: \emph{venīre}, \textbf{veniēbam}; \emph{dormīre}, \textbf{dormiēbat}, he was sleeping.

What it says is the same for all of them: an action under way, or a habit. \emph{Ambulābam} is \textquotedblleft{}I was walking\textquotedblright{}, and also \textquotedblleft{}I used to walk\textquotedblright{}.
\end{grammarlens}

\begin{cousinweb}
Histories of Latin often trace the \textbf{-b-} of \emph{-bam} to the root behind \emph{fuī}, Proto-Indo-European \emph{*b\textsuperscript{h}uH-}, \textquotedblleft{}to become\textquotedblright{}. That is a reconstruction, and neither EtymDB nor Lewis \& Short takes it up, so it is offered as the usual account, not as a fact.

What is certain is where the ending went. Spanish kept it: \emph{cantābat}, \textquotedblleft{}he was singing\textquotedblright{}, is Spanish \emph{cantaba}. French wore the \emph{b} away, so the same form there is \emph{chantait}.
\end{cousinweb}

\subsection*{Your turn}
Say these aloud:

\begin{itemize}
\raggedright
  \item \emph{ambulābam, ambulābās, ambulābat}
  \item \emph{ambulābāmus, ambulābātis, ambulābant}
  \item \emph{habēbam}, I had; \emph{sedēbat}, he was sitting
  \item \emph{Mercātōrēs clāmābant.} — The merchants were shouting
  \item \emph{legēbam}, I was reading; \emph{veniēbam}, I was coming
  \item \emph{Amīcus dormiēbat.} — My friend was sleeping
  \item the closed past, then the open one — \emph{ambulāvit}, then \emph{ambulābat}
\end{itemize}

\begin{tcolorbox}[breakable,colback=teal!4,colframe=teal!35!black,title={Before you move on}]
\textquotedblleft{}They were walking\textquotedblright{}? (\emph{\textbf{Ambulābant.}}) \textquotedblleft{}I had\textquotedblright{}? (\emph{\textbf{Habēbam.}}) \textquotedblleft{}I was reading\textquotedblright{}? (\emph{\textbf{Legēbam.}}) What two letters come before the ending? (\emph{\textbf{-ba-.}}) And which Romance language still keeps the \emph{b}? (\textbf{Spanish}: \emph{cantaba}.)
\end{tcolorbox}
\section[pluēbat … vēnit]{pluēbat … vēnit — the scene and the event}

\label{lesson:LA-C281-pluebat-venit}



\begin{quote}
Say \textquotedblleft{}it rains\textquotedblright{}, with no subject. (\emph{Pluit.}) Now put it in the open past, the way \emph{legere} gives \emph{legēbat}. (\emph{Pluēbat}, it was raining.)
\end{quote}

\begin{grammarlens}[title={Grammar Lens: which past for which job}]
\begin{quote}
\textbf{Pluēbat. Subitō amīcus vēnit.} — It was raining. Suddenly a friend came.
\end{quote}

The rain was already going on, and it went on afterwards: that is the \textbf{imperfect}, \emph{pluēbat}. The friend's arrival happened once, at one moment, and was done: that is the \textbf{perfect}, \emph{vēnit}.

So a story runs on two tracks. The imperfect paints the \textbf{scene}: weather, places, what people were doing. The perfect moves the \textbf{events} along.

\begin{quote}
\textbf{In forō ambulābam; subitō amīcum vīdī.} — I was walking in the forum; suddenly I saw a friend.
\end{quote}

A word like \emph{subitō}, \textquotedblleft{}suddenly\textquotedblright{}, nearly guarantees a perfect after it: suddenness is an event.
\end{grammarlens}

\begin{culture}
Sometimes both are right. \emph{Rōmae habitāvī} counts it up from outside: I lived in Rome, a closed stretch. \emph{Rōmae habitābam} puts you inside it: I was living in Rome then.

One small trap: \emph{pluit} in the perfect is spelled \emph{pluit} too, \textquotedblleft{}it rained\textquotedblright{}, the way \emph{accidit} serves two tenses. Only the story around it decides.
\end{culture}

\begin{cousinweb}
The names are Latin. \textbf{perfectum} comes from \emph{perficiō}, \textquotedblleft{}to carry through, to finish\textquotedblright{} (Lewis \& Short), which EtymDB 2.1 builds from \emph{per-}, \textquotedblleft{}through\textquotedblright{}, and \emph{faciō}, \textquotedblleft{}to make, to do\textquotedblright{}. The perfect is the past that was carried through. \textbf{imperfectum} puts \emph{in-}, \textquotedblleft{}not\textquotedblright{}, in front: unfinished, still going, which is what \emph{pluēbat} is.
\end{cousinweb}

\subsection*{Your turn}
Say these aloud:

\begin{itemize}
\raggedright
  \item \emph{Pluēbat. Subitō amīcus vēnit.}
  \item \emph{In forō ambulābam; subitō amīcum vīdī.}
  \item a scene — \emph{Caelum obscūrum erat.} — and an event — \emph{Amīcus vēnit.}
  \item inside it, then outside it — \emph{Rōmae habitābam}, then \emph{Rōmae habitāvī}
\end{itemize}

\begin{tcolorbox}[breakable,colback=teal!4,colframe=teal!35!black,title={Before you move on}]
Which past paints the scene? (The \textbf{imperfect}: \emph{pluēbat}.) Which moves the story along? (The \textbf{perfect}: \emph{vēnit}.) After \emph{subitō}, which one is likely? (The \textbf{perfect}: suddenness is an event.) And what does \emph{perfectum} mean, taken apart? (\textbf{Carried through}.)
\end{tcolorbox}
\section[(review)]{Review — the scene, and what happened in it}

\label{lesson:LA-R281-the-scene-and-the-event}



\begin{quote}
Say \textquotedblleft{}they saw\textquotedblright{}. (\emph{Vīdērunt.}) Then \textquotedblleft{}we were\textquotedblright{}, the open past. (\emph{Erāmus.})
\end{quote}

\begin{grammarlens}[title={What you've built}]
Four lessons, and both of Latin's story pasts are in your hands. The \textbf{perfect}, a stem and six endings: \emph{vīdī}, \emph{vīdit}, \emph{vīdērunt}, with \emph{-āvī} for most \emph{-āre} verbs. The \textbf{imperfect}: \emph{eram}, \emph{erat}, \emph{erant} for \emph{sum}, and \emph{-ba-} before the ending for every other verb: \emph{ambulābam}, \emph{habēbam}, \emph{legēbam}, \emph{veniēbam}.

And the division of labour between them: the imperfect holds the scene still, the perfect drops the events into it.

\begin{quote}
\emph{Pluēbat, et in caupōnā sedēbam. Subitō amīcus intrāvit.} — It was raining, and I was sitting in the inn. Suddenly a friend came in.
\end{quote}

Two verbs of scene, \emph{pluēbat} and \emph{sedēbam}; one event, \emph{intrāvit}, and \emph{subitō} in front of it.
\end{grammarlens}

\subsection*{Your turn}
Say these aloud:

\begin{itemize}
\raggedright
  \item a scene in three — \emph{Caelum obscūrum erat. Pluēbat. In caupōnā sedēbam.}
  \item the event — \emph{Subitō amīcus intrāvit.} — and a closed one — \emph{In caupōnā cēnāvērunt.}
  \item \emph{In forō ambulābam; subitō amīcum vīdī.}
  \item one of your own — a scene with \emph{erat} or \emph{-bat}, and one thing that happened, with \emph{subitō}
\end{itemize}

\begin{tcolorbox}[breakable,colback=teal!4,colframe=teal!35!black,title={Before you move on}]
\textquotedblleft{}He came in\textquotedblright{}? (\emph{\textbf{Intrāvit.}}) \textquotedblleft{}I was sitting\textquotedblright{}? (\emph{\textbf{Sedēbam.}}) Which ending says \textquotedblleft{}they\textquotedblright{} in the perfect? (\emph{\textbf{-ērunt.}}) In \emph{pluēbat … intrāvit}, which verb is the scene? (\emph{\textbf{Pluēbat.}}) What does \emph{perfectum} mean? (\textbf{Carried through.}) What was the \emph{r} of \emph{eram} once? (An \emph{\textbf{s}}.) And which language still says \emph{cantaba}? (\textbf{Spanish}, keeping the \emph{-b-}.)
\end{tcolorbox}
